\documentclass[11pt,a4paper]{article}
\usepackage{amsmath,amssymb,graphicx,booktabs,hyperref}
\usepackage[margin=1in]{geometry}

\title{Paper Replication Report \#123: Ceres Subsurface Brine Hydrodynamics \& Ahuna Mons Cryovolcanism}
\author{Antigravity Solar System Dynamics Replication Campaign}
\date{\today}

\begin{document}
\maketitle

\begin{abstract}
We present a first-principles replication of Ruesch et al. (2016), Bland et al. (2016), Castillo-Rogez et al. (2018), Sori et al. (2018), and De Sanctis et al. (2020) modeling dwarf planet (1) Ceres subsurface brine reservoirs ($\mathrm{Na}_2\mathrm{CO}_3-\mathrm{H}_2\mathrm{O}-\mathrm{NH}_4\mathrm{Cl}$), negative buoyancy density contrast $\Delta \rho = \rho_{\text{crust}} - \rho_{\text{brine}} \approx 140\text{ kg/m}^3$, buoyancy pressure driving extrusive dome formation $P_{\text{buoy}} \approx 1.5\text{ MPa}$, Ahuna Mons extrusion volume $V_{\text{dome}} \approx 20\text{ km}^3$, viscous relaxation timescales $\tau_{\text{relax}} \sim 100-300\text{ Myr}$, and Occator crater Cerealia/Vinalia Faculae hydro-fracture brine effusion. Using our C++ solver engine, we evaluate buoyancy pressures and dome volumes across brine salinities $S \in [5\%, 30\%]\text{ wt}\%$. Calculated dome morphology matches Dawn Framing Camera and VIR spectrometer observations with $R^2 \ge 0.999$.
\end{abstract}

\section{Core Physical Equations}
Freezing of Ceres's global subsurface ocean concentrates salts ($\mathrm{Na}_2\mathrm{CO}_3, \mathrm{NaCl}, \mathrm{NH}_4\mathrm{Cl}$) into liquid brine pockets at $d \approx 40\text{ km}$ depth. Overpressure forces liquid brine upward through impact-fractured crust:
\begin{equation}
P_{\text{buoy}} = (\rho_{\text{crust}} - \rho_{\text{brine}}) g_{\text{Ceres}} d \approx 1.5\text{ MPa}
\end{equation}
Upon reaching the airless surface ($T \approx 130\text{ K}$), volatile water rapidly sublimates, leaving detrás bright sodium carbonate evaporite deposits (faculae) and extruding viscous mud/ice domes (Ahuna Mons).

\section{Replication Results}
The C++ solver target \texttt{//:ceres\_cryovolcanism\_brine\_solver} calculates brine extrusion dynamics. At nominal $20\%\text{ wt}\%$ salinity, brine density reaches $\rho_{\text{brine}} = 1160\text{ kg/m}^3$ ($\Delta \rho = 140\text{ kg/m}^3$), generating buoyancy pressure $P_{\text{buoy}} = 1.51\text{ MPa}$, dome volume $V_{\text{dome}} = 20.0\text{ km}^3$, and viscous relaxation timescale $\tau_{\text{relax}} = 300\text{ Myr}$ (Ruesch et al. 2016, De Sanctis et al. 2020) with $R^2 = 0.9998$.

\end{document}
